\subsection*{一步一步建立直观}

\paragraph{1. 每个位置都带着一个待演化的量}

常微分方程追踪有限个状态，热方程则追踪每个空间位置的温度。右图的横轴是位置，纵轴可看作相对于平均温度的偏差，负值表示低于平均。时间推进时，整条温度曲线改变，空间形状本身成为系统状态。



\paragraph{2. 温度的弯曲决定当地升温还是降温}

一个温度尖峰比两侧更热，热流向外，峰顶就会下降；谷底则因邻域更热而上升。二阶空间导数把这种邻域不均衡写成局部量。本例的方程是 $u_t=u_{xx}$，且区间首尾相接；边界条件不同，适用的基本模式也会改变。



\paragraph{3. 更细的起伏更快被抹平}

将高频模式的频率加大，而保持观察时间相同。高频部分的曲率更大，所以衰减更快，速度含频率的平方。图形显示的是两个正弦模式的精确演化，不是任意材料与边界的完整仿真；物理单位和扩散系数在这里已经归一化。
